\section{Experimental Setup}
\label{sec:setup}

\paragraph{Models and decoding.}
We evaluate six contemporary LLMs: Gemini~3 Flash (primary ablation model), GPT-5 (low reasoning effort), Claude Haiku~4.5, Grok-4.3, Qwen~3 235B, and DeepSeek Chat. This panel covers a range of closed, open, and general-purpose model families. All six receive the full main evaluation split ($n{=}710$) at baseline. The ledger-feedback diagnostic and other targeted ablations are run on the compact $143$-record ablation split, where every industry--period--difficulty cell and every inconsistency code is represented at least once. Decoding uses temperature $0.0$, except best-of-3 selection (temperature $0.7$, three samples). Additional controls test reasoning level (\appref{appx:reasoning-controls}), provider-native function calling (\appref{appx:native-tools}), character corruption (\appref{appx:ocr-noise}), and run/split stability through repeated decoding and a new seed (\appref{appx:stability-controls}). These diagnostics supplement the six-model headline panel.

\paragraph{Ablations.}
We run four ablation families and report all effects against the same-model baseline. \textbf{(1) Prompt} tests guided solving, self-checking, and stronger citation instructions. Model behavior is comparatively insensitive to these prompt changes. \textbf{(2) Visibility} removes the account list, support documents, metadata, or non-evidence documents. These runs show that accounts and support evidence are necessary and that the evidence-only condition improves document grounding. \textbf{(3) Context stress} adds distractors or moves relevant evidence later to separate longer-context errors from accounting-logic errors. Context hurts performance, while controlled-context records remain difficult. \textbf{(4) Tool/repair} tests calculators, search, the ledger-feedback tool, and best-of-3 sampling. Prompt-described tool access produced no calls, so an additional control exposes the same tools through provider-native function calling. Ledger feedback repairs many reported balance sheets but weakens inconsistency detection. \appref{appx:ablation-grid} gives the evaluated grid, and \appref{appx:native-tools} separates the two tool protocols. For uncertainty, we resample record IDs $5000$ times and recompute each baseline--ablation difference on the same resampled records, yielding intervals for \bsexact{}, \bsrecon{}, strict/lenient joint match, \incCode{}, and document-reference mismatch.

\paragraph{Human validation.}
\label{sec:human-validation}
We released 75 expert-validation records spanning all eight industries, all five difficulty levels, and 15 inconsistency codes. Reviewers were practising finance professionals at two major U.S.\ investment banks and worked from self-contained Markdown records. They assessed accounting plausibility, entry completeness, document support, and calibration, while deterministic generator checks covered arithmetic invariants. Records 1--25 were independently double-reviewed. Collapsing acceptable judgments (``Yes''/``Mostly'' or ``Accept''/``Accept with minor fixes'') gives $100\%$ agreement on document analogy, entries, supporting-document references, and the overall verdict, and $96\%$ on difficulty calibration. Strict exact-option agreement is lower where reviewers differ on severity rather than acceptability: $64\%$ on analogy, $96\%$ on entry correctness, $100\%$ on entry completeness, $72\%$ on citation support, $96\%$ on difficulty, and $80\%$ overall. Separately, a certified accountant independently reviewed the compact coverage split ($143$ records: $120$ standard + $23$ forced-inconsistency). That pass covers every industry--period--difficulty cell and every inconsistency code. The review accepted every record. A small number received minor tightening notes on document references or wording, while all $23$ forced-inconsistency records were confirmed to contain a real contradiction. \appref{appx:human-protocol} gives the protocol, assignment, and marginals.
